
\section{The Ahmad--Cohen scheme}


The computation of the full force for each particle
in the system makes simulations very time--consuming
for large memberships.
Therefore, it is desirable to construct a method in
order to speed up the calculations while retaining
the collisional approach.
One way to achieve this is to employ a ``neighbour scheme'',
suggested by \cite{ac_73}.


The basic idea is to split the force polynomial (\ref{eq:accp0})
on a given particle $i$ into two parts, an irregular and a
regular component:
%
\begin{eqnarray}
\mathbf a_i = \mathbf a_{i,{\rm irr}} + \mathbf a_{i,{\rm reg}} .
\end{eqnarray}
%
The irregular acceleration $\mathbf a_{i,{\rm irr}}$ results
from particles in a certain neighbourhood of $i$
(in the code, FI and FIDOT are the irregular force and its
time derivative at the last irregular step;
internally some routines use FIRR and FD as a local variable).
They give rise to a stronger fluctuating gravitational force,
so it is determined more frequently than the regular one
of the more distant particles that do not change their
relative distance to $i$ so quickly
(in the code, FR and FRDOT are the regular force and its time
derivative at the last regular step; some routines use
as a local variable FREG and FDR).
We can replace the full summation in eq.~(\ref{eq:accel0})
by a sum over the $N_{\rm nb}$ nearest particles for
$\mathbf a_{i,{\rm irr}}$
and add a distant contribution from all the others.
This contribution is updated using another Taylor series
up to the order FRDOT, the time derivative of FR at the
last regular force computation%
%
\footnote{Note, that the code also keeps the variables F and FDOT,
which contain one half (!) of the \emph{total} force,
and one sixth (!) of the \emph{total} time derivative of the force;
this just a handy assignment for the frequent predictions of
equation \ref{eq:predr}. }.

Whether a particle is a neighbour or not is determined
by its distance;
all members inside a specified sphere (``neighbour sphere''
with radius $r_{\rm s}$) are held in a list,
which is modified at the end of each ``regular time--step''
when a total force summation is carried out.
In addition, approaching particles within a surrounding
shell satisfying $\mathbf R \cdot \mathbf V < 0$ are
included.
This ``buffer zone'' serves to identify fast approaching
particles before they penetrate too far inside the
neighbour sphere.
The neighbour criterion should be improved according to relative
forces rather than distances, in particular, if there are very
strong mass differences between particles (black holes!)
--- such kind of work is under progress.


%%%%%%%%%%%%   Figure a    %%%%%%%%%%%%
\noindent
\begin{figure}
\centering
\framebox{
\unitlength0.75cm
\begin{picture}(10,10)
    % star + regularized neighbour + inner circle
    \put(4.9,4.0){$\ast$}
    \put(5.0,4.2){$\diamond$}
    \put(5.0,4.1){\circle{0.6}}
    % neighbour stars + neighbour sphere
    \put(2.3,4.2){$\bullet$}
    \put(4.5,5.3){$\bullet$}
    \put(5.1,6.3){$\bullet$}
    \put(5.8,4.4){$\bullet$}
    \put(6.1,2.2){$\bullet$}
    \put(6.7,2.9){$\bullet$}
    \put(7.6,4.6){$\bullet$}
    \put(7.7,3.6){$\bullet$}
    \put(5.0,4.1){\circle{5.8}}
    % radius of neighbour sphere
    \thinlines
      \put(5.0,4.1){\vector(-1,-2){1.3}}
      \put(4.5,2.5){$r_{\rm s}$}
    % outer stars remaining
    \put(1.7,3.3){$\circ$}
    \put(0.7,5.0){$\circ$}
    \put(0.0,5.3){$\circ$}
    \put(1.5,5.3){$\circ$}
    \put(1.7,5.8){$\circ$}
    \put(0.7,6.6){$\circ$}
    \put(0.6,6.6){$\circ$}
    \put(1.3,7.2){$\circ$}
    \put(3.9,6.8){$\circ$}
    \put(4.3,8.8){$\circ$}
    \put(3.2,9.8){$\circ$}
    \put(4.6,7.2){$\circ$}
    \put(8.7,8.5){$\circ$}
    \put(5.4,8.4){$\circ$}
    \put(5.7,9.5){$\circ$}
    \put(5.8,7.2){$\circ$}
    \put(5.7,7.4){$\circ$}
    \put(7.0,7.9){$\circ$}
    \put(8.0,6.6){$\circ$}
    \put(7.8,6.1){$\circ$}
    \put(9.9,5.0){$\circ$}
    \put(9.1,4.8){$\circ$}
    \put(9.3,3.8){$\circ$}
    \put(3.0,1.7){$\circ$}
    \put(1.1,1.2){$\circ$}
    \put(5.1,0.2){$\circ$}
    \put(7.5,1.8){$\circ$}
    \put(7.7,1.7){$\circ$}
    \put(8.3,1.0){$\circ$}
    \put(8.9,1.0){$\circ$}
\end{picture}
}
   \vspace{1ex}
   \caption{Illustration of the neighbour scheme for
           particle $i$ marked as the asterisk
           (after \cite{a_85}).}
   \label{pic:neighbour}
\end{figure}


Figures \ref{pic:neighbour} and \ref{pic:regtimes} show how
the Ahmad--Cohen scheme works for one particle \cite{ma_92}.
At the beginning of the force calculation, a list of neighbour
objects around the particle $i$ is created first (filled dots).
From this neighbour list the irregular component
$\mathbf a_{i,{\rm irr}}$ is calculated, and then the
summation is continued to the distant particles obtaining
$\mathbf a_{i,{\rm reg}}$.
At the same time we also calculate the first time derivative.
From the equations (\ref{eq:accp0}) and (\ref{eq:accp1})
the position and velocity of the particle $i$ are predicted.
At time $t_{1,{\rm irr}}$ we apply the ``corrector'' only for
$\mathbf a_{i,{\rm irr}}$ from the neighbours;
the regular component we do not correct, but obtain by
extrapolating $\mathbf a_{i,{\rm reg}}$.
At the next step, $t_{2,{\rm irr}}$, the same predictor--corrector
method proceeds for the neighbour particles,
while the correction of the distant acceleration term is
still neglected.
When $t_1$ is reached, the total force is calculated on the
basis of the full application of the Hermite predictor--corrector
method.
Also, a new neighbour list is constructed using the positions
at time $t_1$.   % Makino & Aarseth (1992)
%
Thus, we calculate at certain times only the forces from
neighbours (irregular time--step, $t_{\rm irr}$),
while at other times we calculate both the forces from neighbours
and distant particles (regular time--step, $t_{\rm reg}$).


%%%%%%%%%%%%   Figure b    %%%%%%%%%%%%
\noindent
\begin{figure}[hb]
\centering
\framebox{
\unitlength1cm
\begin{picture}(10,5)
   \thicklines
   % lower horizontal line
      \put(0.5,1){\line(1,0){9.0}}                  % irreg. time
      \multiput(0.5,1.25)(1,0){9}{\line(0,-1){0.5}} % tickmarks
      \put(1.2,0.5){$t_{1,{\rm irr}}$}            % ticknames
      \put(2.2,0.5){$t_{2,{\rm irr}}$}
      \put(3.4,0.5){...}
      \put(5.1,0.5){$t_{1^{\ast},{\rm irr}}$}
      \put(6.1,0.5){$t_{2^{\ast},{\rm irr}}$}
   % upper horizontal line
      \put(0.5,4){\line(1,0){9.0}}                  % regul. time
      \multiput(0.5,4.25)(4,0){3}{\line(0,-1){0.5}} % tickmarks
      \put(0.4,4.5){$t_0$}                       % ticknames
      \put(4.4,4.5){$t_1$}
      \put(8.4,4.5){$t_2$}
   % vertical dotted lines
      \dottedline{0.1}(1.5,1)(1.5,4)
      \dottedline{0.1}(2.5,1)(2.5,4)
      \dottedline{0.1}(3.5,1)(3.5,4)
      \dottedline{0.1}(5.5,1)(5.5,4)
      \dottedline{0.1}(6.5,1)(6.5,4)
   % distance bars
   \thinlines
      \put(0.5,1.6){\vector(1,0){1.0}}
      \put(1.5,1.6){\vector(-1,0){1.0}}
      \put(0.6,1.7){$\Delta t_{\rm irr}$}
      \put(2.0,4.6){\vector(-1,0){1.3}}
      \put(2.1,4.45){$\Delta t_{\rm reg}$}
      \put(3.0,4.6){\vector(1,0){1.3}}
\end{picture}
}
  \vspace{1ex}
  \caption{Regular and irregular time steps (after
           \cite{ma_92}).}
  \label{pic:regtimes}
\end{figure}
%

For a neighbour list of size $N_{\rm nb} \ll N$, this procedure
can lead to a significant gain in efficiency, provided the
respective time scales for $\mathbf a_{i,{\rm irr}}$ and
$\mathbf a_{i,{\rm reg}}$ are well separated.

The actual size of neighbour spheres in NBODY6++ is controlled
iteratively by a requirement in order to keep a certain
optimal number of neighbours.
This variable, NNBOPT, can be adjusted according to performance
requirements.
Its typical values are between 50 and 200 for a very wide range
of total particle numbers $N$.
Outside of the half-mass radius, the requirement of having
NNBOPT neighbours is relaxed due to low local densities.
Insisting on NNBOPT neigbours could result in undesired large
amplitude fluctuations of the neighbour radii.

While \cite{mh_88} claim that the optimal neighbour number
should grow as $N^{3/4}$ (which would be unsuitable for the
performance on parallel computers), this is still an
unsettled question.
\cite{a_85} advocates the coupling of the neighbour radius
to the local density contrast,
%%%
%\begin{eqnarray}
%   \frac{N_{\rm nb}}{N/2}
%          \left( \frac{r_{\rm h}}{r_{\rm s}} \right)^3
%    \approx \frac{n_{\rm nb}}{n_{\rm h}},
%\end{eqnarray}
%%
%where $N$ is the total particle number;
%$r_{\rm h}$ and $r_{\rm s}$ are the half--mass
%radius and the radius of the neighbour sphere;
%and $n_{\rm nb}$ and $n_{\rm h}$ are the
%particle densities inside the neighbour sphere
%and inside the half--mass radius, respectively.
%The approximation in the equation is due to
%the fact, that one--half of the total particle number
%is not necessarily the half of the total mass,
%to which $r_{\rm h}$ is related.
%%%
but NBODY6++ does \emph{not} use that, since it makes average
neighbour numbers much less predictable, which is bad for
the performance and profiling issues on supercomputers, again.

Resuming, the method of the two particle groups is squeezed
into the hierarchical time--step scheme making the overall
view quite complex.
Each particle is moved due to its time--step order \emph{and}
the time--steps, because the force calculation is divided:
In eq.~(\ref{eq:timestep}) a further subscript is needed
which distinguishes the regular and irregular time step.
The accuracy can be tuned by $\eta_{\rm irr} \approx 0.01$
and $\eta_{\rm reg} \approx 0.02$, again.

Both, the neighbour scheme and the hierarchical time--step scheme
have in common that they are centered on one particle $i$,
and they distinguish between nearby and remote stars, and they
save computational time.
One may ask: What is the intriguing difference between them? ---
The neighbour scheme is a \emph{spatial} hierarchy, which
avoids a frequent force calculation of the remote particles,
because their totality provides a smooth potential which
does not vary so much with respect to the particle $i$;
that potential is rather superposed by some fluctuating peaks of
close--by stars which will be ``worked in'' by the more often
force determination.
The time step scheme, in contrast, exhibits the \emph{temporal}
behaviour of the intervals for re--calculation of the full
force in order to maintain the exactness of the trajectory;
time steps chosen too small slow down the advancing calculation
losing the computer's efficiency.
%The neighbour scheme and the individual time--step scheme
%make all calculations non--reversible.
% Funato {\it et al.} (1996) ???
% RSp, Teeseminar 11.10.2001


